% Source: https://artofproblemsolving.com/community/c3h2188094_some_fun_problemweek_1
% Stable identity: aops-topic-2188094-q02
There are $m=\dfrac{n+2}{3}$ terms, and the $k$th term is $(3k-2)(3k-1)(3k)=27k^3-27k^2+6k$. Using $\sum k^3=\left(\dfrac{m(m+1)}{2}\right)^2$, $\sum k^2=\dfrac{m(m+1)(2m+1)}{6}$ and $\sum k=\dfrac{m(m+1)}{2}$ (all from $k=1$ to $m$), the sum is
\[\dfrac{m(m+1)}{4}\left(27m(m+1)-18(2m+1)+12\right)=\dfrac{m(m+1)\left(27m^2-9m-6\right)}{4}=\dfrac{3m(m+1)(3m-2)(3m+1)}{4}.\]
Now $3m-2=n$, $3m+1=n+3$, $m=\dfrac{n+2}{3}$ and $m+1=\dfrac{n+5}{3}$, so the sum equals
\[\boxed{\dfrac{n(n+2)(n+3)(n+5)}{12}}.\]
Check: $n=1$ gives $\dfrac{1\cdot3\cdot4\cdot6}{12}=6=1\cdot2\cdot3$, and $n=4$ gives $\dfrac{4\cdot6\cdot7\cdot9}{12}=126=6+120$.
